\secciongrande{Capacidades y limitaciones de ejecución ¿Por qué javascript en el cliente? Todo sitio web dinámico es gestionado por lenguajes de script en el lado del servidor. En el lado cliente, para mejorar la usabilidad, se recomienda JavaScript por estas razones: 1. No puede escribir archivos en el entorno servidor: Una de sus funcionalidades es preparar los datos que serán enviados, o ceder los datos para que sean tratados en otro lenguaje. 2. Debemos preparar el sitio para que sea navegable sin JavaScript. Lo usamos para crear una navegación rápida, moderna y/o divertida, pero si el navegador no lo tiene habilitado, la página no deberá quedar inactiva. 3. Los datos obtenidos desde el lado del servidor serán examinados con este lenguaje y actualizará las diferentes secciones o partes de la página web.}

\resumen{
  JavaScript se utiliza en el lado cliente para mejorar la usabilidad, permitiendo la preparación de datos, la navegación funcional sin dependencia total de JS, y la actualización dinámica de contenido recibido del servidor. Su principal limitación es la incapacidad de acceder o escribir archivos en el entorno del servidor.
}

\subtituloA{Soluciones que ofrece javascript en el cliente}

\lista{
  **Feedback al usuario:** Dar respuesta inmediata desde formularios o hipertexto. |
  **Interfaz amigable:** Ofrecer una experiencia de usuario moderna y fluida. |
  **Control de elementos:** Gestionar plugins, marcos de navegación o applets de Java. |
  **Procesamiento previo:** Procesar datos de formularios antes de su envío al servidor. |
  **Manipulación dinámica:** Modificar estilos y contenidos de la página en tiempo real dentro del navegador.
}
